% Author: Adam | Module: AI Job Interview Simulator | Project: SmartCareer AI
\section{Introduction et contexte}

SmartCareer AI est une plateforme d'aide à la recherche d'emploi construite en
équipe pour le projet \emph{IA for Software Engineering} 2026-2027. Le
Module~1, réalisé par un coéquipier, génère un CV au format \LaTeX{} à l'aide
d'un petit modèle de langage finement réglé (Qwen3-0.6B + LoRA) servi
localement. Le présent chapitre décrit le \textbf{Module~3 : le simulateur
d'entretien} (\emph{AI Job Interview Simulator}), dont je suis l'unique auteur.

\subsection{Problème}

Un candidat qui dispose d'un CV correct échoue pourtant souvent à l'oral, parce
que personne ne l'a fait répondre avant l'entretien. Les outils d'entraînement
existants reposent presque tous sur une API hébergée : la préparation suppose
donc de confier son CV, son parcours et ses réponses à un tiers. Or un CV est
donnée personnelle au sens du RGPD, et le projet impose déjà un fonctionnement
hors ligne.

\subsection{Objectif}

Le Module~3 propose un entraînement complet, \emph{entièrement local} :

\begin{enumerate}[leftmargin=1.6em,itemsep=1pt,topsep=2pt]
  \item générer des questions d'entretien adaptées à un poste décrit par
        l'utilisateur, éventuellement adossées à un CV déjà produit par le
        Module~1 ;
  \item noter chaque réponse sur 10 selon quatre critères, avec forces,
        faiblesses et verdict ;
  \item proposer un \emph{modèle} de réponse améliorée, structurée en STAR pour
        les questions comportementales ;
  \item conserver l'historique des sessions et une synthèse de clôture.
  \end{enumerate}

Le module s'appuie sur trois leviers, sans aucun entraînement de modèle :
l'ingénierie des prompts, un retrieval-augmented generation (RAG) local sur une
base documentaire écrite pour le projet, et une chaîne de validation stricte de
la sortie du modèle.

\subsection{Contraintes de conception}

Quatre contraintes ont cadré toutes les décisions techniques :

\begin{description}[leftmargin=1.4em,itemsep=2pt,topsep=3pt,style=nextline]
  \item[Coexistence matérielle.] Le Module~1 et le Module~3 partagent le même
    serveur Ollama et la même carte graphique de 4\,096~Mio. Les deux modèles
    ne tiennent pas ensemble sur le GPU : la gestion de la VRAM est une
    fonctionnalité, pas un détail d'exploitation (section~\ref{sec:perf}).
  \item[Non-régression du Module~1.] Le code du coéquipier n'a été touché
    qu'aux endroits strictement nécessaires au branchement : deux lignes dans
    \file{app/\brk{}api.py}, une route dans \file{app.routes.ts}, un lien dans
    \file{workspace.html}, un bloc de variables dans \file{.env.example}
    (section~\ref{sec:branchement}).
  \item[Confidentialité par construction.] Aucune donnée personnelle brute
    n'atteint le modèle, ni les journaux, ni la base (section~\ref{sec:securite}).
  \item[Dégradation contrôlée.] Une panne de modèle, d'index vectoriel ou de
    dépendance lourde ne doit jamais faire échouer un entretien : le module
    retombe sur un repli déterministe explicitement marqué
    \const{degraded} (section~\ref{sec:prompts-validation}).
\end{description}

\subsection{Contributions}

Les livrables sont les suivants.

\begin{itemize}[leftmargin=1.6em,itemsep=1pt,topsep=2pt]
  \item un paquet backend autonome \file{interview\_adam/\brk} : 3\,951 lignes de
        code de production et 2\,368 lignes de tests, soit 6\,319 lignes Python
        au total (mesuré par \file{wc -l}) ;
  \item une interface Angular de 2\,751 lignes (835 lignes TypeScript,
        498 lignes de gabarits HTML, 1\,418 lignes de styles) ;
  \item une base documentaire de 6 domaines, 1\,150 lignes de Markdown et
        71\,495 octets, découpée en 72 fragments indexés ;
  \item 201 tests pytest (200 passés, 1 sauté par défaut) et 12 scénarios
        d'extrémité exécutés dans un navigateur réel ;
  \item un jeu de mesures de latence et d'occupation mémoire, reproduisible
        (section~\ref{sec:perf}).
\end{itemize}

\subsection{Plan}

La section~2 formalise le cahier des charges. La section~3 présente
l'architecture, du diagramme global aux schémas de séquence. La section~4
justifie la pile technique. Les sections~5 à~7 détaillent les trois leviers :
le modèle, les prompts, le RAG. La section~8 traite de sécurité et de
confidentialité, la section~9 des performances mesurées, la section~10 des
tests. La section~11 discute les limites assumées et les perspectives.
Quatre annexes complètent le chapitre : méthode de développement assisté par
IA (A), démonstration (B), environnement d'exécution (C) et contribution
(D).
